\exerciseeditionsection{Solutions to Wald's Problems}
\ExerciseEditorialNote

\begin{enumerate}[leftmargin=2.8em,itemsep=1.2em]
\WaldSolution{1}
\begin{enumerate}[label=\alph*),leftmargin=2em]
\item Concatenation of timelike curves gives
\(I^+(I^+(S))\subseteq I^+(S)\).  Conversely, if \(p\in I^+(S)\), choose a
point \(q\) strictly between the initial point and \(p\) on the timelike
curve.  Then \(q\in I^+(S)\) and \(p\in I^+(q)\), proving equality.

\item Monotonicity gives \(I^+(S)\subseteq I^+(\overline S)\).  If
\(p\in I^+(\overline S)\), some \(q\in\overline S\) lies in the open set
\(I^-(p)\).  That open set meets \(S\), so a point of \(S\) is timelike
related to \(p\).  Hence \(p\in I^+(S)\).
\end{enumerate}

\WaldSolution{2}
\begin{enumerate}[label=\alph*),leftmargin=2em]
\item If \(p\in J^+(S)\), choose points \(p_n\in I^+(p)\) tending to \(p\).
The push-up property gives \(p_n\in I^+(S)\), so
\(p\in\overline{I^+(S)}\).

\item The set \(I^+(S)\) is open and lies in \(J^+(S)\), hence it lies in
the latter's interior.  If \(p\in\operatorname{int}(J^+(S))\), choose
\(q\in I^-(p)\) inside that interior.  Some point of \(S\) is causally
related to \(q\), and the push-up property makes its relation to \(p\)
timelike.  Thus \(p\in I^+(S)\).
\end{enumerate}

\WaldSolution{3}
Suppose first that \(p\in\operatorname{int}(D^+(S))\).  A slightly later
point \(q\) may be chosen in the same interior, so
\(p\in I^-(D^+(S))\); an interior point cannot lie on the achronal initial
set, hence \(p\in I^+(S)\).  Conversely, let
\(p\in I^-(D^+(S))\cap I^+(S)\), and choose \(q\in D^+(S)\cap I^+(p)\).
A small neighborhood \(O\) of \(p\) lies in
\(I^+(S)\cap I^-(q)\).  A past-inextendible causal curve from any point of
\(O\) that missed \(S\) could be joined, using the push-up lemma, to a
past-inextendible causal curve from \(q\) missing \(S\), contradicting
\(q\in D^+(S)\).  Thus \(O\subseteq D^+(S)\).  This proves the first
identity.  Applying it in both time directions and intersecting the two
interior conditions gives
\[
 \operatorname{int}(D(S))
 =I^-(D^+(S))\cap I^+(D^-(S)).
\]

\WaldSolution{4}
Both sets on the right lie in \(J^+(\Sigma)\).  Conversely, take
\(p\in J^+(\Sigma)\).  If \(p\in D^+(C)\), it lies in the first set.  If it
does not, a past-inextendible causal curve through \(p\) avoids \(C\).  Since
\(\Sigma\) is Cauchy, the curve meets \(\Sigma\) at a point of
\(\Sigma-C\).  Therefore \(p\in J^+(\Sigma-C)\).

\WaldSolution{5}
For closed achronal \(S\), Proposition 8.3.2 characterizes
\(\overline{D^+(S)}\) by past-inextendible timelike curves.  Lemma 8.3.3
identifies its interior within \(I^+(S)\).  Removing that interior from the
closure leaves the future horizon.  A point in this remainder lies on
\(\partial I^-(D^+(S))\), and it must lie in \(I^+(S)\cup S\).  Conversely,
a point in
\([I^+(S)\cup S]\cap\partial I^-(D^+(S))\) is in the closure of the domain
but outside its interior.  Hence
\[
 H^+(S)=\bigl[I^+(S)\cup S\bigr]
 \cap\partial I^-(D^+(S)).
\]

\WaldSolution{6}
The time-oriented part of \(\partial D(S)\) lying to the future of \(S\) is
\(\overline{D^+(S)}\setminus\operatorname{int}(D^+(S))=H^+(S)\), using
Lemma 8.3.3 and the definition of the future horizon.  The analogous past
part is \(H^-(S)\).  Their union includes the possible edge points and gives
\[
 \partial D(S)=H^+(S)\cup H^-(S)=H(S).
\]

\WaldSolution{7}
Let \(p_n\in J^+(K)\) converge to \(p\), and choose
\(k_n\in K\) with \(k_n\le p_n\).  Compactness of \(K\) supplies a
subsequence \(k_n\to k\in K\).  Choose \(q\in I^+(p)\); eventually every
causal curve from \(k_n\) to \(p_n\) lies in the compact causal diamond
relevant to \(q\).  The global-hyperbolicity limit-curve theorem yields a
causal limit curve from \(k\) to \(p\).  Thus \(p\in J^+(K)\), proving that
the set is closed.

\WaldSolution{8}
\begin{enumerate}[label=\alph*),leftmargin=2em]
\item Let \(p_n\in J^+(C)\) tend to \(p\), with \(c_n\in C\) and
\(c_n\le p_n\).  Choose \(q\in I^+(p)\).  For large \(n\),
\(c_n\in J^-(q)\cap\Sigma\), which is compact in a globally hyperbolic
spacetime.  A subsequence converges to \(c\in C\), since \(C\) is closed.
The limit-curve theorem gives \(c\le p\), so \(p\in J^+(C)\).

\item In \(2+1\)-dimensional Minkowski spacetime let
\[
 S=\{(-n,-n,1/n):n=1,2,\ldots\}
\]
in coordinates \((t,x,y)\).  This is a closed achronal subset of the null
plane \(t=x\).  Each
\(p_n=(0,0,1/n)\) lies on the future null ray from the corresponding point
of \(S\), and \(p_n\to p=(0,0,0)\).  For every source point, however, the
separation from \(p\) has squared interval \(1/n^2>0\), so it is spacelike.
Thus \(p\notin J^+(S)\), although \(p\in\overline{J^+(S)}\).
\end{enumerate}
\end{enumerate}
